\pretitle{Resumo}
% =========================================================
% RESUMO (elemento obrigatório)
% FINALIDADE: apresentar os pontos principais do trabalho: tema,
%   objetivo, método, resultados e conclusões.
% ESCREVA: em parágrafo único, no lugar do parágrafo de orientação
%   abaixo (depois de \noindent). O título "Resumo" é automático.
% PALAVRAS-CHAVE: troque os exemplos na última linha, separando os termos
%   por ponto e vírgula e terminando com ponto.
% CUIDADO: o abstract deve corresponder a este texto. Extensão e
%   apresentação seguem a ABNT NBR 6028:2021 (de 150 a 500 palavras,
%   voz ativa, terceira pessoa do singular; evite citações, siglas sem
%   definição, fórmulas e remissões a figuras ou tabelas) e as
%   orientações do programa.
% =========================================================
\noindent
Elemento obrigatório. O resumo apresenta, em parágrafo único, os pontos relevantes do trabalho: o tema e o problema investigado, o objetivo da pesquisa, a abordagem metodológica, com indicação dos participantes ou das fontes de dados e dos procedimentos de análise, os principais resultados e as conclusões. Deve ser redigido em frases concisas e afirmativas, com o verbo na voz ativa e na terceira pessoa do singular, e conter de 150 a 500 palavras. Não se recomenda o uso de citações, de siglas sem a forma por extenso, de fórmulas ou de remissões a figuras e tabelas. A primeira frase deve ser significativa e explicar o tema principal do documento. Recomenda-se redigir o resumo após a conclusão do texto, de modo que reflita com fidelidade o conteúdo final da dissertação.

\par\bigskip
% Palavras-chave: termos que representam o conteúdo do trabalho. A
% quantidade segue as orientações do programa e deve ser definida com
% o(a) orientador(a).
\noindent\textbf{Palavras-chave:} palavra-chave 1; palavra-chave 2; palavra-chave 3; palavra-chave 4.
